% Generated draft: requires mathematical and source review.
\documentclass[11pt]{article}
\usepackage[margin=25mm]{geometry}
\usepackage{fontspec}
\setmainfont{Latin Modern Roman}
\usepackage{amsmath,amssymb,mathtools,unicode-math}
\setmathfont{Latin Modern Math}
\usepackage{enumitem,needspace,xurl,hyperref}
\setlength{\parindent}{0pt}
\setlength{\parskip}{5pt}
\setlength{\emergencystretch}{3em}
\newcommand{\N}{\mathbb{N}}
\newcommand{\Z}{\mathbb{Z}}
\newcommand{\Q}{\mathbb{Q}}
\newcommand{\R}{\mathbb{R}}
\newcommand{\C}{\mathbb{C}}
\newcommand{\F}{\mathbb{F}}
\newcommand{\A}{\mathbb{A}}
\newcommand{\PP}{\mathbb{P}}
\newcommand{\E}{\mathbb{E}}
\newcommand{\eps}{\varepsilon}
\newcommand{\dd}{\,\mathrm{d}}
\begin{document}
\section*{UnsolvedMath ID: 30004838}
\textbf{Source problem code:} catalogue entry
\subsection{Definitions and mathematical statement}
The supplied record asks for a characterization of the \(*\)-bialgebras having the L\'evy--Khintchine property \((\mathrm{LK})\). Its terminology requires conventions that are not fully specified in the record.

Under the usual complex, unital, algebraic interpretation, let \(B\) be a unital associative complex \(*\)-algebra with unital \(*\)-homomorphisms
\[
\Delta:B\longrightarrow B\otimes_{\C}B,
\qquad \varepsilon:B\longrightarrow\C,
\]
satisfying coassociativity and the counit identities
\[
(\Delta\otimes\mathrm{id})\Delta=(\mathrm{id}\otimes\Delta)\Delta,
\qquad
(\varepsilon\otimes\mathrm{id})\Delta=\mathrm{id}
=(\mathrm{id}\otimes\varepsilon)\Delta.
\]
Put \(K=\ker\varepsilon\). With these conventions, a generating functional is a complex-linear map \(L:B\to\C\) such that
\[
L(1)=0,\qquad L(a^*)=\overline{L(a)}\quad(a\in B),\qquad
L(a^*a)\geq0\quad(a\in K).
\]
Write \(\mathcal G(B)\) for this class. Under the standard algebraic convention, a generating functional \(G\) is Gaussian when
\[
G(abc)=0\qquad(a,b,c\in K),
\]
equivalently, when \(G\) vanishes on the linear span \(K^3\) of these products. These conventions make explicit a possible intended setting; their agreement with the cited report requires review.

The source also requires a \emph{purely non-Gaussian} generating functional, but supplies no definition of this condition. Let \(\mathcal P(B)\subseteq\mathcal G(B)\) denote the class intended by that source, pending verification of its definition. The requested characterization then has the form: find necessary and sufficient conditions on \(B\) for
\[
\begin{split}
(\mathrm{LK})\qquad
\forall L\in\mathcal G(B)\ \exists G\in\mathcal G(B)\ \exists N\in\mathcal P(B):\quad
&L=G+N,\\
&G(K^3)=\{0\}.
\end{split}
\]
Thus both summands must themselves be generating functionals. No uniqueness is asserted in the supplied statement. The undefined class \(\mathcal P(B)\) is a substantive unresolved qualification: merely requiring that \(N\) fail to be Gaussian does not supply the meaning of ``purely non-Gaussian.'' Consequently, the display records the quantified target but is not a fully specified characterization problem until the source convention for this class and the ambient category are confirmed.
\subsection{Short English statement}
Characterize the \(*\)-bialgebras for which every generating functional is a sum of a Gaussian generating functional and a purely non-Gaussian generating functional. The supplied record does not define the latter condition; the algebraic conventions displayed above also require confirmation against the source. This is a statement-only draft requiring human mathematical review.
\subsection{Sources}
\begin{itemize}
\item Supplied local UnsolvedMath registry record, permanent ID \(30004838\), title ``L\'evy--Khintchine Decomposition for Star-Bialgebras,'' exact source code \texttt{catalogue entry}, in the Oberwolfach Reports Open Problems collection. No proposer is identified; the supplied proposed year is 2021.
\item Supplied authoritative source citation: \emph{Quantum Groups -- Algebra, Analysis and Category Theory} (2022), Oberwolfach Reports, pp.~2397--2458, DOI \texttt{10.4171/owr/2021/44}. The original excerpt identifies p.~2445. Source URL and supplied literature reference ``Oberwolfach Report 44'': \url{https://doi.org/10.4171/owr/2021/44}.
\item Supplied literature reference ``Levy-Khintchine decompositions'': \url{https://arxiv.org/abs/1510.03292}.
\item Supplied literature reference ``Gaussian functionals on easy quantum groups'': \url{https://doi.org/10.1007/s10959-025-01425-w}.
\item The original excerpt attributes the introduction of Gaussian generating functionals in this setting to Sch\"urmann, citing its reference [8]. The bibliographic entry for [8] is not supplied; it cannot be identified here from the record alone.
\item Literature assessment dated exactly \texttt{2026-08-22}: the supplied curation labels the problem ``partially solved'' and reports that property \((\mathrm{LK})\) is proved for commutative bialgebras and several quantum groups, that counterexamples are known, and that no general characterization of \(*\)-bialgebras with this property is available. This is supplied dated curation, not an independently verified literature assessment. The background explicitly calls for independent verification of time-sensitive status and cited claims.
\end{itemize}
\end{document}
